\section{Related work and the novelty claim}
\label{sec:related}

We organise prior work into four bodies and state, for each, what it establishes and what it leaves
open. The section closes with the precise claim this paper defends and an explicit list of claims it
does \emph{not} make.

\subsection{BDD adoption and specification quality}

Empirical work on BDD has concentrated on adoption, scenario structure and maintainability. The
largest study screened 50{,}000 popular open-source projects across five languages, estimated that
about 27\,\% use a BDD framework (68\,\% among Ruby projects), examined co-evolution between scenarios
and production code in 20 Ruby projects using Granger causality --- finding co-evolution in roughly
37\,\% of cases --- and surveyed 31 developers, most of whom reported writing tests while or after
coding rather than strictly applying BDD \citep{zampetti2020demystifying}. Scenario-level analyses have
examined Gherkin structure and smells \citep{chandorkar2022gherkin}, duplicate steps and their
refactoring \citep{binamungu2018duplicate,irshad2022refactoring}, the effort of keeping step
definitions consistent with features \citep{binamungu2018maintaining}, and the characterisation of what
practitioners consider good scenarios \citep{binamungu2020characterising,oliveira2017empirical,
oliveira2019howtoevaluate,couto2023characterization,farooq2023bdd}. Co-evolution between tests and
production code has also been studied outside BDD \citep{Zaidman2008CoEvolution,
Marsavina2014CoEvolution,sun2023coevolution,miranda2025coevolution}.

The limitation, for our purposes, is uniform across this body: quality is assessed by structural
properties, by developer opinion, or by co-change statistics --- never by asking whether the suite
would \emph{fail} if the behaviour changed. A 2023 systematic mapping of 166 BDD papers from 2006 to
2021 makes this concrete: the string ``mutat'' does not occur in the review, and the authors identify
an ``acute shortage of metrics for measuring various aspects of BDD specifications''
\citep{binamungu2023mapping}.

\subsection{Mutation analysis as a measure of discriminating power}

Mutation analysis is the accepted instrument for measuring whether a suite can detect change. The
central validity question --- are mutants a substitute for real faults? --- has been examined
directly: across 357 real faults and roughly 230{,}000 mutants, mutant detection was positively
correlated with real fault detection independently of code coverage, while 17\,\% of faults had no
coupling to any commonly used mutant \citep{Just2014AreMutants}; earlier work questioned whether
mutation is an appropriate experimental tool at all \citep{Andrews2005IsMutation}. Fault revelation
has been studied per criterion \citep{titcheuchekam2017faultrevelation}, redundant and equivalent
mutants are a known measurement hazard \citep{Just2012RedundantMutants,Yao2014EquivalentStubborn,
Kintis2018DetectingTrivial,Kushigian2024EquivalentWild,Straubinger2024ManualEquivalent}, and mutation
testing has been examined as a practitioner activity \citep{Petrovic2021DoesMutationTesting,
Petrovic2022PracticalMutationTesting}. Surveys and threat analyses set out the standards a credible
mutation-based study must meet \citep{Papadakis2019MutationAdvances,Papadakis2016ThreatsValidity,
jia2011survey}.

Two results from this body shape our design directly. First, the coverage--mutation gap: raw
correlations between coverage and mutation score largely disappear once analysis is conditioned on
covered code \citep{Jain2023MindTheGap}, which is why coverage conditioning is mandatory in our
protocol rather than optional. Second, tools are not interchangeable: the same subjects analysed with
four Java tools produced equivalent-mutant counts of 203, 94, 382 and 43, and no tool subsumed the
others \citep{Kintis2018HowEffectiveTools}. Cross-tool pooling of mutation scores is therefore
indefensible, which is why our analysis is stratified by language and tool.

Mutation analysis has been applied at levels above the unit test --- to end-to-end web suites
\citep{leotta2023mutta}, to evolving systems over time \citep{Ojdanic2023MutationEvolving,
ojdanic2023longstanding,akour2020evolution} --- and to specifications themselves
\citep{black2000specmutation}. It has also been studied where practitioners actually are: a large-scale
examination of repositories on GitHub characterises how mutation testing is adopted in the wild
\citep{Sanchez2022MutationInTheWild}, a static and dynamic comparison of Python mutation tools shows
that they diverge in the mutants they generate \citep{Guerino2024PythonTools}, and a study asking
``will my tests tell me if I break this code?'' examines the discriminating power of suites at
different test levels \citep{niedermayr2016willmytests}. The technique has been combined with
test-driven development in a controlled experiment \citep{roman2020tddm} and extended to
cyber-physical systems through data-driven analysis \citep{vigano2022damat}. What it has not been used
for, systematically, is measuring the discriminating power of a \emph{BDD specification layer} across
a corpus.

\subsection{The oracle problem and assertion quality}

Whether a test can detect a fault depends on its oracle \citep{barr2015oracleproblem}, and the choice of
oracle has been evaluated empirically, including the use of assertions specifically
\citep{shrestha2011assertionsoracles}. Assertion quality has been studied as a defect source
\citep{zhi2013inadequateassertions,zhang2015assertions}, as a measurable density
\citep{Catolino2019AssertionDensity}, and as a driver of false confidence in ``pseudo-tested'' methods
that are covered but not effectively asserted \citep{veraperez2018pseudotested}. Test-suite assessment
has looked at whether suites would catch injected faults
\citep{yandrapally2021maewu,soto2021e2edataset}. Static indicators have been used to
\emph{predict} mutation scores at the level of an individual test case
\citep{Grano2021Lightweight}, and removing redundant tests has been shown to change coverage and
mutation scores \citep{Bhatta2025TestDeletion}. At the system level, smells specific to system user
interactive tests have been catalogued \citep{rwemalika2023suit}, and test-code refactoring has been
related to test-code quality outcomes \citep{Martins2024TestRefactoring}.

The nearest miss to our RQ3 is \citet{Grano2021Lightweight}: static indicators estimating a test
case's mutation score. Our unit is different (the repository's specification layer, not an individual
test case), our indicators are drift-oriented rather than complexity-oriented, and our outcome is
BDD-only kill rather than overall mutation score. We cite it as the closest precedent, not as a
displaced one.

\subsection{Decay, drift and erosion}

The vocabulary of decay is established for architecture and design, not for test suites.
Architectural drift and erosion were characterised early \citep{PerryWolf1992}, with subsequent work on
architectural drift detection \citep{Rosik2010ArchDrift} and architectural evolution
\citep{Wermelinger2011ArchEvolution}; code decay and design-pattern decay have likewise been
formalised \citep{Izurieta2012Decay}. On the test side, the relevant evidence concerns disabled tests
\citep{Kim2021Disabled}, test-case survival \citep{Feldt2014GrowOld}, rotten green tests
\citep{Delplanque2019RottenGreen,Aranega2021RottenGreen}, test-smell prevalence and survival
\citep{Bavota2015AreTestSmells,Bavota2012Distribution,kim2020testsmell,Kim2021SecretLife,
Spadini2018Relation,Spadini2020Severity,Pontillo2024ML,Panichella2022TwentyYears,
Cai2026PairedSmells,Martins2023Diffusion,Fushihara2023Trend,Moonen2008Interplay,
Athanasiou2014TestCodeQuality,Pinto2012Myths}, and co-change between feature files and code
\citep{yang2019cochanges}.

Two terminological points must be made plainly, because the literature does not license the terms we
would otherwise like to use. \emph{Specification drift} is \textbf{not} an established term in
software engineering: our searches of Crossref and OpenAlex found no peer-reviewed source that
establishes it. The same holds for \emph{test decay}, \emph{test erosion} and \emph{test rot} as
applied to test suites. We therefore introduce \emph{specification drift} explicitly as a definition
--- a specialisation of architectural drift to the executable-specification layer --- and we do not
present it as a received construct. This is a terminological contribution at best, and we do not
overstate it.

\subsection{Traceability in the Information Systems literature}\label{sec:traceability}

This subsection is \emph{motivating background}, not a report of a measurement. Requirements
traceability between a requirement and the artefacts that implement and verify it is a long-standing
IS concern, motivated by the need to demonstrate that a system satisfies the requirements it was built
against; a systematic review of pre-requirements-specification traceability characterises its
motivations and the persistent difficulty of establishing it in practice
\citep{mucha2024prerequirements}. That difficulty is not confined to plan-driven settings: Brazilian
studies of agile practice report that requirements traceability remains problematic when requirements
are expressed as user stories and acceptance criteria rather than as numbered specifications
\citep{trindade2016rastreabilidade}, and that independently evolving systems in a software ecosystem
complicate tracing further \citep{goncalves2022rastreabilidade}.

Two consequences follow for this paper, and we separate them carefully because the distinction was
blurred in an earlier revision. First, traceability is the reason the construct matters
organisationally: a scenario that references no identifiable business rule may still be a good test,
but it cannot connect the system's behaviour to a stated obligation, and the IS literature explains why
such connections are processes to be maintained rather than states to be achieved
\citep{henderson1993strategic}. Second, \textbf{this paper does not measure traceability coverage.} An
earlier version of the protocol defined such a measure --- the fraction of identifiable business rules
that at least one scenario references --- but identifying business rules reliably requires human coding
of repository documentation, and no such coding was performed here. The measure is therefore defined in
the protocol and left unimplemented; no traceability figure appears anywhere in this paper, and none
should be inferred from its presence in the protocol. We note it as an explicit extension rather than
presenting a number we did not collect.

\subsection{What is new, precisely}

The claim this paper defends is a conjunction of four elements, each individually unclaimed:

\begin{quote}
No study has measured, on a corpus of public Cucumber-style suites, the mutation-based oracle strength
of the BDD layer \emph{alone} --- separated from the contribution of the repository's other tests ---
nor how that strength changes across the repository's history.
\end{quote}

The closest prior work is a single industrial proof of concept that applied mutation analysis to a BDD
acceptance suite for a robotic system and reported that 79\,\% of generated mutants were detected, but
used hand-selected domain mutants, a single system, no public corpus, no separation of the BDD layer's
contribution from other tests, and no temporal analysis \citep{santos2024aat4irs}. An industrial
experience report on mutation testing in a Ruby project is sometimes cited as BDD evidence, but its
published artefacts concern Ruby mutation operators and practitioner adoption; the only openly
available version does not mention Cucumber or Gherkin at all \citep{li2015mutationpracticeruby}. We
therefore do not describe it as a BDD study. Model-based generation of Cucumber scenarios
\citep{li2016skyfire} addresses a different question.

Accordingly, we explicitly do \textbf{not} claim that: BDD quality is unstudied; BDD suites have never
been analysed at scale; mutation testing has never been applied to a BDD suite; longitudinal mutation
analysis is new; or that nobody has related static test-quality indicators to mutation scores. Each of
these would be false, and the corresponding prior work is cited above.
